\section{相关工作}

\subsection{受限理性元推理与智能体过度思考}
机载算力分配源于受限理性与元推理理论，核心是评估“计算的价值”是否超过其时间与能量代价\cite{horvitz1987reasoning,russell1991right}。在机器人运动规划中，研究者通过学习停止规则或自适应计算预算来决定何时终止搜索并执行动作，或使规划与移动并行交织\cite{xiao2020thinking,sung2021quit,elboher2023formal,magnino2026metareasoning}。然而，当内部环境模型存在偏差时，延长规划视野或增加搜索深度反而会放大模型误差，诱发决策分析中的“优化器诅咒”，使规划器陷入无益的过度思考\cite{smith2006curse,jiang2015horizon,cuadron2025overthinking}。近期在智能体推理与任务消歧中也发现了相似规律：当外部环境状态未知或指令存在多重解释时，增加测试时计算无法替代与环境交互或主动寻求澄清\cite{hunt2021formalizing,shen2025thinking,kobalczyk2025active}。与仅关注何时停止搜索的元推理不同，本文将模糊性拆解为求解器粗糙引起的数值歧义与地图失真引起的环境歧义，从而在计算、观测与保守行动之间建立切换准则。

\subsection{面向决策的信息获取与多解性判别}
并非所有减少环境不确定性的观测都能改善特定决策，唯有能区分竞争假设的信息才具决策价值\cite{chernoff1959sequential,box1967discrimination}。在统计实验设计与多臂赌博机中，最优臂识别与信息导向采样专门收集能尽快排除次优选项的样本，而非均匀降低全局方差\cite{garivier2016optimal,russo2018thompson,russo2018ids}；目标导向实验设计同样强调只针对目标泛函优化传感器布放\cite{attia2018goal}。与之呼应，信息间隙决策理论与预测多解性研究指出，当多个与数据同样相容的模型给出相反预测时，单靠点估计无法保证安全，必须考察不确定性半径内的最坏表现\cite{benhaim2006infogap,marx2020predictive}。不同于在离散假设空间或线性反问题中计算期望信息增益，本文利用包络定理直接在连续无散海流场上沿梯度搜寻翻转二元航路选择的物理反例，以最坏后悔度量环境歧义的严重程度。

\subsection{不确定海洋流场中的航路规划与自适应采样}
强海流下的自主水下航行器路径规划常利用图搜索或水平集方程在确定性流场中求取时间最优轨迹\cite{dijkstra1959note,lolla2014levelset}。为应对海洋预报不确定性，研究者提出了集合期望与风险折中规划、动态正交随机时间最优控制以及纳入误差协方差的鲁棒路径搜索\cite{wang2016path,subramani2018stochastic,kularatne2018forecast}；海洋自适应采样则利用高斯过程或集合卡尔曼滤波在线更新叶绿素或流场估计，引导航行器前往高信息增益区\cite{anderson2001eakf,srinivas2012gpucb,hollinger2014sampling,fossum2019}。然而，现有不确定流场规划大多将集合样本作为固定输入来寻找单一路径，既不评估有限集合是否漏掉了翻转拓扑走向的流场，也不回答是否值得付出时间代价去换取新海流图；本文则将多年海洋再分析大数据转化为合理流场集合的物理边界，专门解决粗图下“算、测、走”的闭环元决策问题。
